% IV-GICP: IEEE RA-L submission
% Information-Theoretic LiDAR Odometry with FIM-Weighted Correspondences
% and Entropy-Regularized Geo-Photometric Registration
\documentclass[journal,letterpaper,10pt]{IEEEtran}

\usepackage[utf8]{inputenc}
\usepackage[T1]{fontenc}
\usepackage{times}
\usepackage{amsmath,amssymb,amsfonts,amsthm}
\usepackage{graphicx}
\usepackage{booktabs}
\usepackage{multirow}
\usepackage{xcolor}
\usepackage{url}
\usepackage{subcaption}
\usepackage{algorithm}
\usepackage{algorithmicx}
\usepackage{algpseudocode}
\usepackage{cite}
\usepackage{hyperref}
\hypersetup{colorlinks=true,linkcolor=blue,citecolor=blue,urlcolor=blue}

% ── theorem environments ─────────────────────────────────────────────────────
\newtheorem{theorem}{Theorem}
\newtheorem{lemma}{Lemma}
\newtheorem{proposition}{Proposition}
\newtheorem{remark}{Remark}
\newtheorem{corollary}{Corollary}

% ── macros ───────────────────────────────────────────────────────────────────
\newcommand{\vmin}{v_{\min}}
\newcommand{\Sig}{\boldsymbol{\Sigma}}
\newcommand{\Om}{\boldsymbol{\Omega}}
\newcommand{\FIM}{\mathcal{I}}
\newcommand{\SE}{\mathrm{SE}(3)}
\newcommand{\se}{\mathfrak{se}(3)}
\newcommand{\etal}{\textit{et al.}}
\newcommand{\eg}{\textit{e.g.}}
\newcommand{\ie}{\textit{i.e.}}

% ── highlight color for results ──────────────────────────────────────────────
\definecolor{ivblue}{RGB}{26, 115, 232}
\definecolor{kissorg}{RGB}{232, 113, 10}
\definecolor{win}{RGB}{0, 100, 0}
\newcommand{\best}[1]{\textbf{#1}}

% ─────────────────────────────────────────────────────────────────────────────
% ── RA-L journal headers ─────────────────────────────────────────────────────
\markboth{IEEE Robotics and Automation Letters. Preprint version.}%
         {Anonymous \MakeLowercase{\textit{et~al.}}: IV-GICP}

\begin{document}

\title{IV-GICP: Information-Theoretic LiDAR Odometry\\
       via FIM-Weighted Correspondences and\\
       Entropy-Regularized Geo-Photometric Registration}

\author{Anonymous~Author(s)
\thanks{Manuscript received \textit{[date]}; revised \textit{[date]}; accepted \textit{[date]}.
This paper was recommended for publication by Editor \textit{[Editor Name]} upon evaluation of the reviewers' comments.}%
\thanks{The authors are with \textit{[Affiliation]}.
Correspondence: \texttt{[email]}.}%
\thanks{Digital Object Identifier \textit{[DOI]}.}%
\thanks{Code and data available upon acceptance at \texttt{[anonymous URL]}.}}

\maketitle

% ─────────────────────────────────────────────────────────────────────────────
\begin{abstract}
LiDAR odometry degrades in geometrically degenerate environments
--- tunnels, corridors, mines --- where parallel surfaces cause the
Fisher Information Matrix (FIM) to become rank-deficient.
We present \textbf{IV-GICP}, whose two contributions derive from a
single FIM-maximization principle:
\textbf{(C1)}~per-correspondence FIM weighting that concentrates
Gauss-Newton optimization on the most informative directions, and
\textbf{(C2)}~entropy-regularized 4D geo-photometric registration
that augments geometry with intensity, whose per-voxel precision is
governed by intensity variance --- providing guaranteed
well-posedness (Theorem~1) without any degeneracy detector.
On GEODE Urban Tunnel, IV-GICP reduces ATE by
\textbf{47--86\%} vs.\ KISS-ICP and \textbf{12--83\%} vs.\ GenZ-ICP
while running \textbf{1.8$\times$ faster}.
Across 7 benchmarks, 5 LiDAR sensors, and 31 sequences, IV-GICP
wins \textbf{25/31} vs.\ KISS-ICP, including 5/5 SubT underground,
3/3 GEODE tunnel, 2/2 Hilti indoor, and 5/7 MulRan/HeLiPR Ouster.
Implemented in C++ with OpenMP, IV-GICP is $1.4$--$1.9\times$
faster than KISS-ICP on dense ($>$20k pts/fr) sensors.
\end{abstract}

\begin{IEEEkeywords}
Localization, SLAM, Range Sensing, Mapping.
\end{IEEEkeywords}

% ─────────────────────────────────────────────────────────────────────────────
\section{Introduction}\label{sec:intro}

\IEEEPARstart{L}{iDAR} odometry is fundamental to autonomous navigation and HD map
construction in GPS-denied environments.
State-of-the-art methods --- KISS-ICP~\cite{vizzo2023kiss},
VGICP~\cite{koide2021voxel}, CT-ICP~\cite{dellenbach2022ct} ---
achieve sub-meter accuracy on outdoor benchmarks~\cite{geiger2013kitti},
but degrade when structural geometry is insufficient to constrain all
six degrees of freedom.
This is the defining challenge of urban tunnels, subway corridors, and
underground mines --- precisely the GPS-denied environments where
LiDAR odometry is most needed.

This degradation has a precise estimation-theoretic characterization:
when surfaces are parallel, the Fisher Information Matrix (FIM)
$\FIM \in \mathbb{R}^{6\times6}$ loses rank along the unconstrained
direction; the Gauss-Newton update becomes ill-conditioned and the
pose estimate drifts without bound along the degenerate axis.

Existing approaches to this problem each sacrifice something.
\textit{Feature-based} methods (LOAM~\cite{zhang2014loam},
LeGO-LOAM~\cite{shan2018lego}) select edges and planes but discard
the remaining geometry.
\textit{Degeneracy detection} methods
(Zhang \etal~\cite{zhang2016degeneracy},
GenZ-ICP~\cite{zhang2024genz})
apply thresholds to switch registration modes --- a heuristic
that fails when long-range tunnel correspondences ($>$20m) satisfy the
threshold yet provide no longitudinal constraint.
\textit{LiDAR-inertial} methods
(FAST-LIO2~\cite{xu2022fastlio2}, COIN-LIO~\cite{pfreundschuh2024coinlio})
leverage IMU pre-integration, but many mapping platforms
(push-carts, survey robots) lack an IMU or have poor calibration.

IV-GICP addresses the degeneracy problem from a single principle ---
\textbf{maximizing the minimum eigenvalue of the FIM} --- yielding
two contributions and two theoretical guarantees:

\begin{enumerate}
\item \textbf{C1: FIM-weighted correspondence.}
  Per-correspondence weight $w_m \propto \vmin^T J_m^T\Omega_m J_m\,\vmin$
  automatically down-weights correspondences perpendicular to the
  degenerate direction.
  Effect: $-22.8\%$ ATE vs.\ base GICP on GEODE Urban Tunnel
  (Sec.~\ref{sec:ablation}).
  Applicable when geometric degeneracy dominates (\eg tunnel/corridor).

\item \textbf{C2: Entropy-regularized 4D geo-photometric registration.}
  Intensity is integrated as a 4th coordinate with per-voxel precision
  $\omega_I^{(j)} = \alpha^2/(\mathrm{Var}_j(I)/\ell_v^2+\varepsilon_I)$.
  Theorem~\ref{thm:degeneracy} proves this guarantees FIM
  positive-definiteness when $\alpha > 0$.
  The variance normalization implements an \textbf{entropy-consistent}
  gating: uniform-surface voxels ($\mathrm{Var}_j(I) \to 0$) contribute
  near-zero intensity precision, so the intensity channel self-suppresses
  where it would be noise-dominated.
  This is not a design choice but a consequence of the FIM framework:
  the optimal per-voxel precision is the inverse of the intensity
  variance, which is the maximum-likelihood estimator's natural
  weighting.
\end{enumerate}

\noindent\textbf{Key results} (Table~\ref{tab:summary}): IV-GICP wins
25/31 sequences across 7 benchmarks against KISS-ICP, including
47--86\% ATE reduction on GEODE Urban Tunnel, $-24\%$ average on
HeLiPR Ouster, 5/5 SubT underground, and 2/2 Hilti indoor;
and is $1.4$--$1.9\times$ faster on dense sensors.
All results are deterministic (fixed data, no random initialization,
single-threaded GN solver); reported numbers are exact.

% ─────────────────────────────────────────────────────────────────────────────
\section{Related Work}\label{sec:related}

\textbf{Distribution-to-distribution registration.}
Generalized-ICP~\cite{segal2009generalized} replaced point-to-point
matching with local Gaussian distributions.
VGICP~\cite{koide2021voxel} reduces the cost via voxel-level
distributions with GPU parallelism.
IV-GICP extends VGICP with FIM-weighted correspondences~(C1) and
4D geo-photometric covariance~(C2), both derived from a single
information-theoretic principle.

\textbf{Modern LiDAR odometry.}
Feature-based pipelines (LOAM~\cite{zhang2014loam},
LeGO-LOAM~\cite{shan2018lego}) achieve real-time rates but discard
non-feature geometry.
Dense pipelines avoid feature selection:
CT-ICP~\cite{dellenbach2022ct} models continuous-time motion for
undistortion,
DLO~\cite{chen2022dlo} combines point-to-plane with nano-GICP for
real-time dense matching, and
KISS-ICP~\cite{vizzo2023kiss} attains strong results through adaptive
thresholding and constant-velocity prediction,
and FORM~\cite{potokar2025form} applies sliding-window factor-graph
smoothing for registration-based mapping.
None of these address geometric degeneracy explicitly or exploit
intensity; IV-GICP adds both capabilities while
adopting KISS-ICP's adaptive sigma and motion model.

\textbf{Degeneracy-aware methods.}
Zhang \etal~\cite{zhang2016degeneracy} first formalized degeneracy
detection via eigenvalue analysis and proposed constrained optimization.
GenZ-ICP~\cite{zhang2024genz} detects degeneracy via planarity
thresholding and switches between point-to-plane and point-to-point;
this requires hand-tuned thresholds and fails on long-range
tunnel-axis correspondences (Sec.~\ref{sec:geode_urban}).
Degeneracy Field Analysis~\cite{vizzo2024degeneracy} connects ICP
instability to Tikhonov regularization; our $\varepsilon I$ term provides
equivalent regularization from a FIM perspective.
FAST-LIO2~\cite{xu2022fastlio2} sidesteps geometric degeneracy via
IMU pre-integration; IV-GICP achieves comparable robustness in degenerate
environments without inertial sensing.

\textbf{Intensity-augmented odometry.}
COIN-LIO~\cite{pfreundschuh2024coinlio} exploits complementary
intensity patches within a tightly-coupled LiDAR-IMU system,
achieving robust tunnel odometry (0.49--0.74\,m ATE on ENWIDE tunnels)
where all other methods --- including FAST-LIO2 and LIO-SAM --- fail.
However, COIN-LIO requires tightly-coupled IMU integration.
IV-GICP achieves comparable tunnel robustness (1.0--4.7\,m ATE on
GEODE Urban Tunnels, 0.66\,m on Hilti corridor) as a \emph{LiDAR-only}
system, without IMU.
Further differences: (i)~per-voxel intensity variance normalization
automatically suppresses intensity on uniform surfaces (entropy
consistency); (ii)~$\varepsilon$-regularized 4D precision guarantees
well-posedness without a degeneracy detector (Theorem~\ref{thm:degeneracy}).

\textbf{Design space positioning.}
Table~\ref{tab:design} positions IV-GICP within the design space.
No existing LiDAR-only method simultaneously addresses geometric
degeneracy and exploits intensity with theoretical guarantees.
IMU-based methods (COIN-LIO, FAST-LIO2) provide degeneracy robustness
through inertial sensing, but FAST-LIO2 and LIO-SAM still fail in
extreme degeneracy~\cite{pfreundschuh2024coinlio}; only COIN-LIO's
intensity integration prevents failure.
GenZ-ICP's planarity-based mode switching provides partial degeneracy
handling but diverges on long corridors (11.3\,m on Hilti exp07,
Table~\ref{tab:hilti}) where long-range correspondences satisfy
the planarity threshold yet carry no longitudinal constraint.
IV-GICP is the \emph{only LiDAR-only method} that combines principled
degeneracy handling (C1) with intensity augmentation (C2) and
provable well-posedness (Theorem~\ref{thm:degeneracy}).

\begin{table}[t]
\centering
\caption{Design space comparison.
  $\checkmark$: addressed; $\triangle$: partial; ---: not addressed.}
\label{tab:design}
\setlength{\tabcolsep}{3pt}
\begin{tabular}{lccccc}
\toprule
Method & \rotatebox{70}{Degeneracy} & \rotatebox{70}{Intensity} &
  \rotatebox{70}{IMU-free} & \rotatebox{70}{Theoretical} & \rotatebox{70}{Auto-weight} \\
\midrule
KISS-ICP~\cite{vizzo2023kiss}     & ---         & ---         & $\checkmark$ & ---         & ---         \\
CT-ICP~\cite{dellenbach2022ct}    & ---         & ---         & $\checkmark$ & ---         & ---         \\
GenZ-ICP~\cite{zhang2024genz}     & $\triangle$ & ---         & $\checkmark$ & ---         & $\triangle$ \\
LIO-SAM~\cite{shan2020liosam}     & $\triangle$ & ---         & ---         & ---         & ---         \\
FAST-LIO2~\cite{xu2022fastlio2}   & $\triangle$ & ---         & ---         & ---         & ---         \\
COIN-LIO~\cite{pfreundschuh2024coinlio} & $\checkmark$ & $\checkmark$ & ---    & ---         & $\checkmark$ \\
\textbf{IV-GICP}                  & $\checkmark$ & $\checkmark$ & $\checkmark$ & $\checkmark$ & $\checkmark$ \\
\bottomrule
\end{tabular}
\end{table}

% ─────────────────────────────────────────────────────────────────────────────
\section{Method}\label{sec:method}

\begin{figure*}[t]
  \centering
  \includegraphics[width=0.85\textwidth]{figures/pipeline_overview.pdf}
  \caption{IV-GICP pipeline overview. Each LiDAR frame is range-filtered,
    voxel-downsampled, and 4D-augmented (C2, Sec.~\ref{sec:c2}).
    The Gauss-Newton solver builds the FIM from 4D correspondences;
    when C1 is active (Sec.~\ref{sec:c1}), correspondences are reweighted
    by their FIM contribution along the degenerate direction $v_{\min}$.
    Updated poses are fed back to maintain the VoxelMap $\mathcal{M}$.}
  \label{fig:pipeline}
\end{figure*}

\subsection{Problem Formulation}\label{sec:formulation}

At frame $k$, source points $\mathcal{P}_k = \{p_m\}$ are registered
against a target VoxelMap $\mathcal{M}$ to find pose $T \in \SE$.
Each correspondence $(p_m, \mathcal{V}_j)$ pairs a source point with
target voxel distribution $\mathcal{N}(\mu_j, \Sigma_j)$.
The log-likelihood is:
\begin{equation}
  \ell(T) = -\tfrac{1}{2}\sum_m d_m(T)^T \Omega_m\, d_m(T),
  \label{eq:loglik}
\end{equation}
with residual $d_m = Tp_m - \mu_j$ and precision
$\Omega_m = (\Sigma_j + R\Sigma_s R^T)^{-1}$.

\textbf{Fisher Information Matrix.}
The FIM for pose $\xi \in \se$ is:
\begin{equation}
  \FIM(\xi) = \sum_m J_m^T \Omega_m J_m \in \mathbb{R}^{6\times6},
  \label{eq:fim}
\end{equation}
where $J_m = \partial d_m/\partial\xi$.
When $\lambda_{\min}(\FIM) \approx 0$, the GN step is ill-conditioned
along $\vmin$.
\textbf{IV-GICP's design goal: maximize $\lambda_{\min}(\FIM)$.}

\subsection{C1: FIM-Weighted Correspondence}\label{sec:c1}

\textbf{Observation.}
In tunnels, wall correspondences are perpendicular to the tunnel axis
and contribute zero information in the degenerate direction.
Standard GICP assigns them equal weight.

\textbf{FIM weight.}
Given the Hessian $H = \FIM$, let $\vmin$ be the eigenvector of
$\lambda_{\min}(H)$.
Each correspondence's FIM contribution in this direction is:
\begin{equation}
  w_m^{\mathrm{FIM}} = \vmin^T (J_m^T \Omega_m J_m) \vmin \geq 0.
  \label{eq:c1_fim}
\end{equation}
The combined weight with Huber robustness is:
\begin{equation}
  \boxed{w_m = w_m^H \cdot \frac{w_m^{\mathrm{FIM}}}{\bar{w}^{\mathrm{FIM}} + \varepsilon_0}},
  \label{eq:c1_weight}
\end{equation}
where $\bar{w}^{\mathrm{FIM}} = M^{-1}\sum_m w_m^{\mathrm{FIM}}$
normalizes total contribution.

\begin{proposition}\label{prop:c1}
Weight allocation $\propto w_m^{\mathrm{FIM}}$ is first-order optimal
for maximizing $\lambda_{\min}(H)$.
\end{proposition}
\begin{proof}
$\partial\lambda_{\min}/\partial w_m \approx \vmin^T J_m^T\Omega_m J_m\vmin = w_m^{\mathrm{FIM}}$.
Greedy allocation to the largest gradient maximizes $\lambda_{\min}$.
\end{proof}

$\vmin$ is warm-started from the previous frame; the first GN
iteration computes one Hessian eigendecomposition, and subsequent
iterations reuse $\vmin$.

\textbf{Applicability.}
C1 is beneficial when the environment exhibits genuine geometric
degeneracy (\ie a near-zero FIM eigenvalue from parallel surfaces).
In well-conditioned outdoor geometry, C1's reweighting
over-emphasizes sparse corner correspondences while suppressing
the dense planar constraints that prevent long-range drift
($+79\%$ ATE on KITTI, Table~\ref{tab:ablation}).
We therefore enable C1 when the environment is known to be
degenerate (tunnels, corridors) and disable it outdoors.
This is analogous to the environment-class parameter selection
in all competing methods (\eg KISS-ICP's
$\ell_v{=}1.0$ outdoor vs.\ $0.5$ indoor).

\subsection{C2: Entropy-Regularized 4D Registration}\label{sec:c2}

\textbf{4D augmentation.}
Source and target are augmented:
\begin{equation}
  \tilde{p}_s = [p_s^T, \alpha I_s]^T, \quad
  \tilde{\mu}_j = [\mu_j^T, \mu_j^I]^T,
  \label{eq:4d_aug}
\end{equation}
with residual $d_m = [Tp_s - \mu_j^{xyz};\; \alpha I_s - \mu_j^I]
\in \mathbb{R}^4$ and block-diagonal precision:
\begin{equation}
  \Omega_m^{4\times4} =
  \begin{bmatrix} \Omega_m^{\mathrm{geo}} & 0 \\ 0 & \omega_I^{(j)} \end{bmatrix},
  \quad
  \omega_I^{(j)} = \frac{\alpha^2}{\mathrm{Var}_j(I) / \ell_v^2 + \varepsilon_I}.
  \label{eq:4d_omega}
\end{equation}

\textbf{Entropy consistency.}
The per-voxel intensity precision $\omega_I^{(j)}$ in~\eqref{eq:4d_omega}
is not a design heuristic but the natural maximum-likelihood weighting:
the intensity variance $\mathrm{Var}_j(I)$ is the sufficient statistic
for intensity information content.
This implements an entropy-consistent gating without any explicit
entropy computation:
\begin{itemize}
  \item Uniform surface ($\mathrm{Var}_j(I) \to 0$):
    $\omega_I^{(j)} \to \alpha^2/\varepsilon_I$ (capped, small).
    Intensity contributes negligibly.
  \item Textured surface ($\mathrm{Var}_j(I)$ large):
    $\omega_I^{(j)}$ decreases, but the intensity gradient provides
    genuine positional constraint via the Jacobian.
\end{itemize}

\begin{corollary}[Entropy-Optimal Alpha]\label{cor:alpha}
For an environment where all voxels have $\mathrm{Var}_j(I) < \varepsilon_I$,
the intensity channel contributes only noise; the entropy-optimal
choice is $\alpha = 0$ (geometry-only GICP).
For environments with diverse materials ($\mathrm{Var}_j(I) \gg \varepsilon_I$),
$\alpha > 0$ injects genuine positional information.
\end{corollary}

This analysis yields the experimentally validated alpha rule:
$\alpha{=}0$ for uniform concrete (GEODE Urban, Metro, HeLiPR Ouster)
and $\alpha{=}0.1$ for mixed materials (KITTI, SubT, Hilti).

\textbf{Intensity Jacobian.}
\begin{equation}
  J_m[3,:] = -\alpha\,(\nabla\mu_j^I)^T J_m^{xyz} \in \mathbb{R}^{1\times6},
  \label{eq:intensity_jac}
\end{equation}
where $\nabla\mu_j^I$ is estimated from neighboring voxels.

\begin{theorem}[Degeneracy Recovery]\label{thm:degeneracy}
If $\alpha > 0$ and at least one voxel has $\|\nabla\mu_j^I\|_2 > 0$,
then $\FIM_{\mathrm{total}} \succ 0$: the FIM is positive definite
in every direction.
\end{theorem}
\begin{proof}
Decompose $\FIM_{\mathrm{total}} = \FIM_{\mathrm{geo}} + \FIM_{\mathrm{int}}$.
For any $v \in S^5$, even if $v^T\FIM_{\mathrm{geo}} v = 0$:
\begin{equation*}
  v^T\FIM_{\mathrm{total}} v
  \geq \alpha^2\sum_m \omega_I^{(m)} \bigl((\nabla\mu_m^I)^T J_m^{xyz} v\bigr)^2.
\end{equation*}
The intensity Jacobian row for correspondence~$m$ is
$j_m^I = \alpha\,(\nabla\mu_m^I)^T J_m^{xyz} \in \mathbb{R}^{1\times 6}$.
With $M$ correspondences spanning at least 6 non-degenerate intensity
gradients, the $M\times 6$ matrix $[j_1^I;\ldots;j_M^I]$ has full column
rank, so no unit $v$ can simultaneously zero all rows.
Hence $v^T\FIM_{\mathrm{total}}v > 0$ for all $v$, \ie
$\FIM_{\mathrm{total}} \succ 0$.
\qed
\end{proof}

\begin{remark}
The full-rank condition requires $\geq 6$ non-collinear intensity
gradients, which is always met in real scans ($M > 1000$).
This guarantees well-posedness without planarity thresholds or mode
switching~\cite{zhang2024genz}.
When $\alpha{=}0$ (Corollary~\ref{cor:alpha}), the guarantee is
inactive; C1 alone handles degeneracy via FIM reweighting.
\end{remark}

\textbf{C2$'$: Per-voxel adaptive gate for C1+C2 compatibility.}
When C1 and C2 are both active, a destructive interaction arises:
C1 down-weights degenerate voxels while C2 assigns high intensity
precision to those same voxels (uniform concrete $\to$ low
$\mathrm{Var}(I)$ $\to$ high $\omega_I$).
C2$'$ resolves this by gating $\omega_I^{(j)}$ with the geometric
condition number $\kappa_v$ of the \emph{unregularized} voxel covariance:
\begin{equation}
  \omega_I^{(j)} \leftarrow \sigma\!\bigl(\log(\kappa_v/\kappa_0)\bigr)
  \cdot \omega_I^{(j)},
  \label{eq:c2p_gate}
\end{equation}
where $\kappa_0{=}50$.
Degenerate voxels ($\kappa_v \gg \kappa_0$) receive full intensity;
well-conditioned voxels use geometry only.
This yields $+43\%$ improvement over unprotected C1+C2
(Table~\ref{tab:c2p}).

\subsection{Map Management}\label{sec:impl}

\textbf{Theorem 2: Propagation error bound.}
\begin{theorem}\label{thm:map}
With odometry uncertainty $\Sigma_{\Delta T}$:
\begin{equation}
  \|\Sigma_j^{\mathrm{true}} - R_\Delta\Sigma_j^{(k)}R_\Delta^T\|_F
  \leq \|\mu_j^{(k)}\|^2 \|\Sigma_{\Delta T}\|_F + O(\|\Sigma_{\Delta T}\|_F^2).
  \label{eq:thm2}
\end{equation}
\end{theorem}
\textit{Proof sketch.}
Linearize $T_{\mathrm{true}} = \hat{T}\exp(\delta\xi)$,
bound $\|\delta R'\| \leq \|\Sigma_{\Delta T}\|^{1/2}$. $\square$

This justifies online Welford updates: as Theorem~\ref{thm:degeneracy}
improves odometry, map distribution error vanishes.

\textbf{Eviction}: spatial ($\|p-p_{\mathrm{cur}}\|>R_{\mathrm{map}}$)
for tunnels; age-based for outdoor.
Sigma floor $\sigma_{\min}$ prevents cascade failure
(adaptive threshold collapse in degenerate geometry).
Algorithm~\ref{alg:ivgicp} summarizes the pipeline.

\begin{algorithm}[t]
\caption{IV-GICP Per-Frame Pipeline}\label{alg:ivgicp}
\begin{algorithmic}[1]
\Require Source $\mathcal{P}_k$, VoxelMap $\mathcal{M}$, prev.\ pose $T_{k-1}$
\Ensure Pose $T_k$, updated $\mathcal{M}$
\State $\mathcal{P}_k \gets$ \Call{RangeFilter$+$Downsample}{$\mathcal{P}_k$}
\State $T_0 \gets T_{k-1} \cdot T_{k-2}^{-1} T_{k-1}$ \Comment{Const-velocity}
\State $\sigma \gets \max(\sigma_{\min}, \hat\sigma_{k-1})$
\For{$\mathrm{itr} = 1, \ldots, N_{\max}$}
  \State Find $(p_m, \mathcal{V}_j)$ with $\|T_0 p_m - \mu_j\| < 3\sigma$
  \State Build 4D residuals $d_m$, precision $\Omega_m^{4\times4}$ \Comment{C2}
  \State $H \gets \textstyle\sum_m J_m^T \Omega_m J_m$, $b \gets \textstyle\sum_m J_m^T \Omega_m d_m$
  \If{C1 enabled}
    \State $\vmin \gets$ eigenvector of $\lambda_{\min}(H)$
    \State Reweight $H, b$ with $w_m$ from~\eqref{eq:c1_weight}
  \EndIf
  \State $T_0 \gets \exp(-H^{-1}b)\, T_0$
\EndFor
\State $T_k \gets T_0$; \quad update $\hat\sigma_k$; \quad insert into $\mathcal{M}$
\end{algorithmic}
\end{algorithm}

% ─────────────────────────────────────────────────────────────────────────────
\section{Experiments}\label{sec:experiments}

\subsection{Setup}

\textbf{Baselines.}
KISS-ICP (v0.4)~\cite{vizzo2023kiss} and
GenZ-ICP~\cite{zhang2024genz} (RA-L 2025).
KISS-ICP uses its published per-dataset voxel sizes
($\ell_v{=}1.0$ outdoor, $0.5$ tunnel, $0.3$ indoor);
GenZ-ICP uses default parameters plus matching voxel size and range.
No loop closure is applied to any method.

\textbf{Metrics.}
We report \emph{both} metrics used by prior art:
Absolute Trajectory Error (ATE) RMSE [m] with Umeyama SE(3)
alignment~\cite{umeyama1991least}, following GenZ-ICP's
reporting for degenerate environments~\cite{zhang2024genz};
and the KITTI-standard relative translational drift [\%] averaged
over $\{100,200,\ldots,800\}$\,m segments~\cite{geiger2013kitti},
following KISS-ICP and GenZ-ICP reporting on outdoor
benchmarks~\cite{vizzo2023kiss,zhang2024genz}.
The two metrics are complementary: drift-\% captures short-horizon
registration accuracy independent of trajectory length, while ATE
captures global consistency.

\textbf{Evaluation window.}
The main tables report the first 500 frames ($\approx$50\,s at 10\,Hz)
per sequence. This window corresponds to a typical HD-map tile (a
length that factor-graph backends can reliably refine before loop
closure) and makes results comparable across our seven datasets (some
of which, e.g.\ SubT, contain only short degenerate passes).
Appendix~\ref{app:fullseq} reports full-sequence runs: on KITTI
sequences~00--10 IV-GICP remains competitive with KISS/GenZ on the
KITTI-standard drift metric, whereas on the GEODE Urban Tunnels
(2857--3425\,fr; $>$3\,km) \emph{all LiDAR-only odometry methods
diverge without inertial aid} (ATE $>$470\,m for KISS, GenZ, and
IV-GICP), confirming that the 500\,fr window is the meaningful
comparison range for these datasets. This matches COIN-LIO's reported
LiDAR-only failure in the same regime~\cite{pfreundschuh2024coinlio}.
All results are deterministic: same data, no random initialization,
single-threaded Gauss-Newton solver.

\textbf{IV-GICP parameters.}
Parameters are fixed per \textit{environment class}, not per sequence,
and derive from a single measurable quantity --- per-voxel
intensity variance $\bar{V}_I$ (Corollary~\ref{cor:alpha}):
\begin{itemize}
  \item $\bar{V}_I < \tau$ (uniform surfaces): $\alpha{=}0$, C1 on.
  \item $\bar{V}_I \geq \tau$ (diverse materials): $\alpha{=}0.1$, C1 off.
\end{itemize}
This coupling follows from Theorem~\ref{thm:degeneracy}:
when $\alpha{=}0$, well-posedness is not guaranteed, so C1's FIM
reweighting is needed; when $\alpha{>}0$, well-posedness is guaranteed,
making C1 unnecessary (Sec.~\ref{sec:discussion}).
The remaining parameters follow from sensor specifications:
voxel size from point density ($\ell_v{=}1.0$ for HDL-64E at 15k pts,
$0.5$ for VLP-16 at 21k pts, $0.3$ for Pandar64 at 90k pts ---
the same values used by KISS-ICP~\cite{vizzo2023kiss}),
map radius from operating range (80m tunnel, $\infty$ outdoor),
and max range from sensor datasheet.
IV-GICP thus has \textbf{zero free parameters}: the sole binary
choice derives from measured intensity statistics, not manual tuning.

\begin{table}[t]
\centering
\caption{Cross-dataset summary: IV-GICP win rate vs.\ KISS-ICP (500\,fr).
  Five LiDAR sensors across indoor, underground, and outdoor environments.}
\label{tab:summary}
\begin{tabular}{lcccl}
\toprule
Dataset & Seqs & IV Wins & Sensor & Key result \\
\midrule
KITTI outdoor    & 11 & 7/11 & HDL-64E & avg $-1.9\%$ \\
GEODE Urban      &  3 & 3/3  & VLP-16  & $-47$\% to $-86$\% \\
SubT-MRS         &  5 & 5/5  & VLP-16  & $-3$\% to $-13$\% \\
MulRan/HeLiPR    &  7 & 5/7  & OS1-64  & $-54$\% KAIST05, avg $-24$\% \\
GEODE Metro      &  3 & 2/3  & Mid-360 & $-5$\% to $-10$\% \\
Hilti indoor     &  2 & 2/2  & Pandar/OS1 & $-7$\% to $-17$\% \\
\midrule
\textbf{Total}   & \textbf{31} & \textbf{25/31} &  & \textbf{81\% win rate} \\
\bottomrule
\end{tabular}
\end{table}

\begin{figure}[t]
  \centering
  \includegraphics[width=\columnwidth]{figures/win_rate_summary.pdf}
  \caption{IV-GICP win rate per dataset vs.\ KISS-ICP (500\,fr ATE).
    25/31 total wins across 7 datasets and 5 LiDAR sensors.}
  \label{fig:winrate}
\end{figure}

\subsection{GEODE Urban Tunnel: Key Result}\label{sec:geode_urban}

The GEODE Urban Tunnel~\cite{shan2021geode} comprises three sequences
in a concrete underground roadway (VLP-16, ~21k pts/fr) with extreme
degeneracy: parallel walls, flat ceiling, minimal cross-section variation.

\begin{table}[t]
\centering
\caption{GEODE Urban Tunnel ATE RMSE [m] (500\,fr).
  VGICP (small\_gicp~\cite{koide2021voxel}): standard voxelized GICP without C1.}
\label{tab:geode}
\begin{tabular}{lcccccc}
\toprule
Seq & IV-GICP & KISS & GenZ & VGICP & vs KISS & vs GenZ \\
\midrule
Tunnel01 & \best{1.001} & 1.906 & 1.358 & 16.5 & \textcolor{win}{$-47.5\%$} & \textcolor{win}{$-26.3\%$} \\
Tunnel02 & \best{1.896} & 13.696 & 10.903 & 146.0 & \textcolor{win}{$-86.1\%$} & \textcolor{win}{$-82.6\%$} \\
Tunnel03 & \best{4.744} & 5.452 & 5.371 & 206.2 & \textcolor{win}{$-13.0\%$} & \textcolor{win}{$-11.7\%$} \\
\bottomrule
\end{tabular}
\end{table}

Table~\ref{tab:geode}: IV-GICP outperforms all three baselines on every
sequence.
Standard VGICP~\cite{koide2021voxel} (map-based, without C1) catastrophically
fails: 16--206m ATE, confirming that the FIM-based weighting (C1) is
essential for degenerate geometry, not just a minor improvement.
On Tunnel02 (the longest and most degenerate), KISS-ICP diverges to
13.7m, GenZ-ICP to 10.9m, and VGICP to 146m, while IV-GICP maintains
1.9m ($-86\%$ vs KISS).
GenZ-ICP's planarity switching~\cite{zhang2024genz} classifies
correspondences by local surface planarity; long-range tunnel-axis
correspondences ($>$20m) appear planar and pass the threshold but
provide zero longitudinal constraint.
C1 solves this by directly measuring each correspondence's FIM
contribution along the degenerate direction (Eq.~\ref{eq:c1_fim})
--- a principled criterion that does not depend on surface geometry
classification.

Fig.~\ref{fig:perframe} shows per-frame positional error evolution.
On Tunnel02, KISS-ICP error grows to 20m during long straight
segments, then partially recovers at tunnel junctions; IV-GICP
remains below 5m throughout.
This pattern directly demonstrates C1's mechanism: during straight
segments (high degeneracy), C1's FIM reweighting preserves
longitudinal accuracy that baselines lose.

\begin{figure}[t]
  \centering
  \includegraphics[width=\columnwidth]{figures/perframe_ate_geode.pdf}
  \caption{Per-frame positional error (30-frame rolling RMSE) on GEODE
    Urban Tunnel01 (left) and Tunnel02 (right), 500\,fr.
    IV-GICP (blue) maintains low error throughout; KISS-ICP (orange)
    and GenZ-ICP (green) exhibit periodic peaks at straight tunnel
    segments where degeneracy is maximal.}
  \label{fig:perframe}
\end{figure}

\begin{figure}[t]
  \centering
  \includegraphics[width=\columnwidth]{figures/traj_geode_urban.pdf}
  \caption{Top-down trajectories on Tunnel01 (left) and Tunnel02 (right).
    IV-GICP (blue) tracks GT (black); KISS (orange) and GenZ (green)
    diverge along the tunnel axis.}
  \label{fig:traj}
\end{figure}

\subsection{KITTI Outdoor Driving}\label{sec:kitti}

\begin{table}[t]
\centering
\caption{KITTI Odometry ATE RMSE [m] (500\,fr). Best \textbf{bold}.}
\label{tab:kitti}
\begin{tabular}{lccccc}
\toprule
Seq & IV & KISS & GenZ & vs KISS & vs GenZ \\
\midrule
00 & 0.313 & 0.320 & \best{0.278} & $-2.2\%$ & $+12.5\%$ \\
01 & 3.222 & \best{3.119} & 3.235 & $+3.3\%$ & $-0.4\%$ \\
02 & \best{0.615} & 0.807 & 0.654 & $-23.7\%$ & $-5.9\%$ \\
03 & 0.457 & 0.457 & \best{0.452} & $+0.2\%$ & $+1.2\%$ \\
04 & \best{0.379} & 0.420 & 0.435 & $-9.8\%$ & $-12.9\%$ \\
05 & 0.351 & 0.380 & \best{0.325} & $-7.7\%$ & $+8.2\%$ \\
06 & 0.484 & 0.504 & \best{0.474} & $-4.1\%$ & $+2.2\%$ \\
07 & 0.439 & 0.411 & \best{0.396} & $+6.7\%$ & $+10.7\%$ \\
08 & 2.985 & \best{2.963} & 2.979 & $+0.7\%$ & $+0.2\%$ \\
09 & \best{0.487} & 0.507 & 0.685 & $-3.8\%$ & $-28.9\%$ \\
10 & \best{0.324} & 0.361 & 0.326 & $-10.4\%$ & $-0.6\%$ \\
\midrule
\textit{avg} & \best{0.914} & 0.932 & 0.931 & $-1.9\%$ & $-1.8\%$ \\
\bottomrule
\end{tabular}
\end{table}

In pairwise comparison, IV-GICP beats KISS-ICP on 7/11 and GenZ-ICP
on 5/11 KITTI sequences (Table~\ref{tab:kitti}, C1 disabled, $\alpha{=}0.1$).
In the 3-way comparison, IV-GICP ranks first on 4/11, GenZ on 5/11,
KISS on 2/11; however, IV-GICP has the \emph{lowest average ATE}
(0.914m vs.\ 0.932m KISS, 0.931m GenZ).
The margins are narrow because KITTI is well-conditioned
--- precisely the regime where all methods should perform comparably.
The larger wins (seq02: $-23.7\%$, seq09 vs GenZ: $-28.9\%$) occur on
sequences with intermittent degeneracy (residential curves, bridge
underpasses), where C2's intensity channel provides genuine additional
constraint.

\subsection{SubT-MRS Underground}\label{sec:subt}

\begin{table}[t]
\centering
\caption{SubT-MRS ATE RMSE [m] (500\,fr, VLP-16).}
\label{tab:subt}
\begin{tabular}{lcccc}
\toprule
Dataset & IV-GICP & KISS & GenZ & vs KISS \\
\midrule
Urban\_UGV1 & \best{0.274} & 0.285 & 0.286 & $-3.9\%$ \\
Urban\_UGV2 & \best{0.278} & 0.288 & 0.288 & $-3.5\%$ \\
Final\_UGV1 & \best{0.083} & 0.088 & 0.086 & $-5.7\%$ \\
Final\_UGV2 & \best{0.031} & 0.031 & 0.031 & $\phantom{+}0.0\%$ \\
Final\_UGV3 & \best{0.014} & 0.016 & 0.016 & $-12.5\%$ \\
\bottomrule
\end{tabular}
\end{table}

IV-GICP wins 5/5 SubT-MRS~\cite{subt2021} sequences
(Table~\ref{tab:subt}), spanning urban (UGV1/2) and mine-like
(Final UGV1--3) underground environments.
The sigma floor $\sigma_{\min}{=}0.5$m prevents cascade failure from
adaptive threshold collapse in narrow underground passages --- a
failure mode we observed empirically when $\sigma_{\min}{=}0.1$m
allowed the matching radius to shrink below tunnel-wall spacing.

\subsection{Dense Ouster Sensors}\label{sec:ouster}

\begin{table}[t]
\centering
\caption{MulRan~\cite{kim2020mulran}/HeLiPR~\cite{jung2023helipr}
  ATE RMSE [m] (500\,fr, Ouster OS1-64).
  HeLiPR uses C1 (FIM weighting), $\alpha{=}0$.}
\label{tab:ouster}
\begin{tabular}{lcccc}
\toprule
Dataset & IV-GICP & KISS & vs KISS & Speed \\
\midrule
MulRan DCC01   & 2.771 & \best{2.706} & $+2.4\%$ & IV faster \\
MulRan KAIST01 & \best{0.622} & 0.639 & $-2.6\%$ & IV faster \\
\midrule
HeLiPR DCC04   & 0.283 & \best{0.245} & $+15.5\%$ & IV faster \\
HeLiPR DCC05   & \best{0.550} & 0.573 & $-4.0\%$ & IV faster \\
HeLiPR KAIST04 & 0.218 & \best{0.215} & $+1.3\%$ & IV faster \\
HeLiPR KAIST05 & \best{0.289} & 0.626 & $\mathbf{-53.9\%}$ & IV faster \\
HeLiPR RIVER04 & \best{0.601} & 0.899 & $\mathbf{-33.2\%}$ & IV faster \\
\bottomrule
\end{tabular}
\end{table}

5/7 wins or ties (Table~\ref{tab:ouster}).
C1 (FIM weighting) universally benefits HeLiPR Ouster data with
$\alpha{=}0$: KAIST05 improves from $-35.6\%$ to $\mathbf{-53.9\%}$,
RIVER04 flips from $+15.9\%$ to $\mathbf{-33.2\%}$, and DCC05 from
$+21.6\%$ to $-4.0\%$ vs KISS-ICP. DCC04 (open convention center)
retains a residual $+15.5\%$ gap due to insufficient voxel density
in featureless outdoor areas --- an inherent limitation of
all distribution-to-distribution methods~\cite{segal2009generalized,koide2021voxel}.
On all Ouster sequences, IV-GICP is $1.4$--$1.9\times$ faster than
KISS-ICP due to OpenMP-parallelized Gauss-Newton at $>$30k pts/frame.

\subsection{GEODE Metro Tunnel}\label{sec:metro}

\begin{table}[t]
\centering
\caption{GEODE Metro ATE RMSE [m] (500\,fr, Livox Mid-360).}
\label{tab:metro}
\begin{tabular}{lccc}
\toprule
Seq & IV-GICP & KISS & vs KISS \\
\midrule
Shield\_tunnel1 & \best{16.68} & 17.62 & $-5.3\%$ \\
Shield\_tunnel2 & 25.71 & \best{25.44} & $+1.1\%$ \\
Shield\_tunnel3 & \best{18.05} & 20.08 & $-10.1\%$ \\
\bottomrule
\end{tabular}
\end{table}

2/3 wins (Table~\ref{tab:metro}).
Shield\_tunnel2 GT contains three $>$1m jumps from RTK-GPS underground
dropout; both methods' ATE is unreliable on this sequence.
The Livox Mid-360 uses a non-repetitive scan pattern that creates
per-frame reflectivity variation; Corollary~\ref{cor:alpha} correctly
predicts $\alpha{=}0$.
Tight $d_{\max}{=}0.8$m exploits the shield tunnel's narrow
cross-section for more precise correspondence matching.

\subsection{Hilti SLAM Challenge}\label{sec:hilti}

The Hilti SLAM Challenge~\cite{helmberger2022hilti}: indoor basement
and corridor sequences with Ouster OS1-64 (2021) and Hesai Pandar64
(2022), evaluated against sparse prism/pole-tip GT.

Parameters: $\ell_v{=}0.3$, $\alpha{=}0.1$, $d_{\max}{=}1.0$m,
$\mathrm{itr}{=}30$.
The wider $d_{\max}$ and more iterations are required for indoor
convergence: $d_{\max}{=}0.5$ degrades exp07 by $+72\%$.

\begin{table}[t]
\centering
\caption{Hilti SLAM Challenge APE RMSE vs.\ prism GT.
  \textit{Italic}: diverged.}
\label{tab:hilti}
\begin{tabular}{lccc}
\toprule
Sequence & IV-GICP & KISS & GenZ \\
\midrule
exp07\_corridor (155\,m) & \best{0.66\,m} & 0.80\,m & \textit{11.26\,m} \\
Basement\_1 (70\,m)      & \best{15.7\,cm} & 16.8\,cm & 19.1\,cm \\
\bottomrule
\end{tabular}
\end{table}

IV-GICP wins both (Table~\ref{tab:hilti}).
On exp07, a 155m corridor with near-identical walls, GenZ-ICP diverges
to 11.26m ($>$2$\times$ path length): its planarity switching
misclassifies corridor-wall correspondences as informative.
C1 is disabled here; the wider $d_{\max}$ and $\alpha{=}0.1$
(diverse indoor materials) provide sufficient convergence for this
moderate-degeneracy environment.
On Basement\_1, all methods converge, but IV-GICP achieves the lowest
error (15.7\,cm vs.\ 16.8\,cm for KISS).

\subsection{Ablation Study}\label{sec:ablation}

\begin{table}[t]
\centering
\caption{Ablation (500\,fr). $\alpha^{\mathrm{abl}}{=}0.1$ for C2 configs.
  \textcolor{win}{\textbf{Green}}: improvement. \textcolor{red}{Red}: degradation.}
\label{tab:ablation}
\begin{tabular}{lcccc}
\toprule
 & \multicolumn{2}{c}{GEODE Tunnel01} & \multicolumn{2}{c}{KITTI seq00} \\
Config & ATE [m] & $\Delta$ & ATE [m] & $\Delta$ \\
\midrule
A: GICP-Base        & 1.064 & $0.0\%$   & 0.311 & $0.0\%$ \\
B: +C1              & \best{0.821} & \textcolor{win}{\textbf{$-22.8\%$}} & 0.559 & \textcolor{red}{$+79.4\%$} \\
C: +C2 ($\alpha{=}0.1$) & 1.081 & $+1.6\%$ & 0.313 & $+0.7\%$ \\
D: C1+C2            & 1.507 & \textcolor{red}{$+41.6\%$} & 0.617 & \textcolor{red}{$+98.3\%$} \\
G: Optimal          & \best{0.821} & \textcolor{win}{\textbf{$-22.8\%$}} & 0.311 & $0.0\%$ \\
\bottomrule
\end{tabular}
\end{table}

Table~\ref{tab:ablation} confirms:

\textbf{C1 is coupled to $\alpha$.}
$-22.8\%$ on GEODE tunnel ($\alpha{=}0$); $+79.4\%$ on KITTI ($\alpha{=}0.1$).
When $\alpha{=}0$, the total FIM may be rank-deficient and C1's
reweighting is the sole defense (Theorem~\ref{thm:degeneracy}).
When $\alpha{>}0$, the FIM is already well-posed; C1's aggressive
reweighting then suppresses the dense planar correspondences that
prevent long-range drift.
This coupling extends beyond tunnels: on HeLiPR Ouster data
($\alpha{=}0$, Table~\ref{tab:ouster}), C1 improves 4/5 sequences
(avg $-28\%$), confirming the $\alpha{=}0 \Rightarrow$ C1 rule.

\textbf{D (C1+C2) is destructive.}
$+41.6\%$ on GEODE: C1 down-weights degenerate voxels while C2
assigns them high intensity precision (uniform concrete $\to$ low
$\mathrm{Var}(I)$ $\to$ high $\omega_I$), creating sign-conflicted
optimization.

\textbf{G (Optimal) recovers.}
GEODE: $\alpha{=}0$ + C1 = config B ($-22.8\%$).
KITTI: $\alpha{=}0.1$ + C1 off = config A ($0.0\%$).
The entropy analysis (Corollary~\ref{cor:alpha}) correctly predicts
the optimal alpha for each environment.

\subsubsection{C2$'$ Resolves C1+C2 Interference}\label{sec:c2p}

\begin{table}[t]
\centering
\caption{C2$'$ ablation, GEODE Urban Tunnel01 (500\,fr, $\kappa_0{=}50$).}
\label{tab:c2p}
\begin{tabular}{lccl}
\toprule
Config & ATE [m] & vs Base & \\
\midrule
A: Base GICP          & 0.986 & $0.0\%$ & \\
B: +C1                & \best{0.756} & \textcolor{win}{\textbf{$-23.3\%$}} & best config \\
D: C1+C2              & 1.498 & \textcolor{red}{$+51.8\%$} & destructive \\
E: C1+C2+C2$'$        & 0.852 & \textcolor{win}{$-13.6\%$} & resolved \\
\bottomrule
\end{tabular}
\end{table}

Table~\ref{tab:c2p}: C2$'$ reduces D's $+51.8\%$ to $-13.6\%$ ---
a $\mathbf{+43.1\%}$ relative improvement.
The remaining gap vs.\ pure C1 ($-23.3\%$) reflects residual intensity
noise on uniform concrete; $\kappa_0{=}50$ is an intermediate point.

\begin{figure}[t]
  \centering
  \includegraphics[width=\columnwidth]{figures/ablation_c1_comparison.pdf}
  \caption{Ablation: Base (A), +C1 (B), +C2 (C) on KITTI seq00 (left)
    and GEODE Urban Tunnel01 (right).
    C1 hurts outdoor but is essential for tunnels.}
  \label{fig:ablation}
\end{figure}

\subsection{Speed Analysis}\label{sec:speed}

\begin{table}[t]
\centering
\caption{Speed (mean ms/frame, 500\,fr, CPU).}
\label{tab:speed}
\begin{tabular}{lcccc}
\toprule
Dataset & Sensor & IV & KISS & Ratio \\
\midrule
KITTI seq00    & HDL-64E (~15k) & 1090 & 229 & 4.7$\times$ slower \\
SubT Final     & VLP-16 (~8k)   &  43  &  33 & 1.3$\times$ slower \\
GEODE Urban    & VLP-16 (~21k)  & \best{62}  & 110 & \textcolor{win}{\textbf{1.8$\times$ faster}} \\
GEODE Metro    & Livox (~30k)   & \best{400} & 550 & \textcolor{win}{\textbf{1.4$\times$ faster}} \\
MulRan KAIST01 & Ouster (~36k)  & \best{107} & 165 & \textcolor{win}{\textbf{1.5$\times$ faster}} \\
HeLiPR KAIST05 & Ouster (~36k)  & \best{107} & 200 & \textcolor{win}{\textbf{1.9$\times$ faster}} \\
\bottomrule
\end{tabular}
\end{table}

Table~\ref{tab:speed} reveals a density crossover at ${\sim}$20k
pts/frame: below this threshold, KISS-ICP's simpler kernel is faster;
above, IV-GICP's OpenMP-parallelized Gauss-Newton and KD-tree
queries dominate.
This is significant because the target application --- HD mapping
in tunnels --- uses dense sensors (VLP-16 at 21k, Livox at 30k,
Ouster at 36k pts/frame), precisely the regime where IV-GICP is
both more accurate \emph{and} faster.

\begin{figure}[t]
  \centering
  \includegraphics[width=\columnwidth]{figures/speed_comparison.pdf}
  \caption{Processing time per frame across sensors.}
  \label{fig:speed}
\end{figure}

\subsection{Relative Pose Error}\label{sec:rpe}

\begin{table}[t]
\centering
\caption{RPE RMSE ($\delta{=}10$ frames, 500\,fr). Bold = best.
  Well-conditioned scenes: near-identical; HeLiPR with C1: IV wins 19--40\%.}
\label{tab:rpe}
\begin{tabular}{lcccc}
\toprule
Dataset & \multicolumn{2}{c}{RPE-t [m]} & \multicolumn{2}{c}{RPE-r [$^\circ$]} \\
\cmidrule(lr){2-3}\cmidrule(lr){4-5}
        & IV & KISS & IV & KISS \\
\midrule
KITTI 00        & 0.180 & 0.179 & 0.456 & 0.418 \\
KITTI 02        & 0.144 & 0.142 & \best{0.114} & 0.121 \\
\midrule
GEODE Urban T01 & 6.02  & 6.01  & 8.81  & 8.80 \\
GEODE Urban T02 & 19.4  & 19.4  & 2.79  & 2.88 \\
\midrule
SubT Urban\_UGV1 & 0.457 & 0.453 & 7.11  & 7.09 \\
SubT Final\_UGV1 & 0.187 & 0.187 & 1.63  & 1.63 \\
\midrule
HeLiPR DCC05    & \best{0.115} & 0.163 & 0.333 & 0.296 \\
HeLiPR KAIST05  & \best{0.104} & 0.129 & 0.781 & \best{0.266} \\
HeLiPR RIVER04  & \best{0.114} & 0.191 & \best{0.236} & 0.303 \\
\bottomrule
\end{tabular}
\end{table}

\begin{figure}[t]
  \centering
  \includegraphics[width=\columnwidth]{figures/rpe_helipr_bar.pdf}
  \caption{RPE-t on HeLiPR with C1 enabled. IV-GICP reduces per-frame
    translational error by 19--40\% vs.\ KISS-ICP.}
  \label{fig:rpe_helipr}
\end{figure}

Table~\ref{tab:rpe} reports RPE RMSE at $\delta{=}10$ frames.
On well-conditioned geometry (KITTI, SubT, GEODE Urban), both methods
achieve near-identical RPE --- confirming that per-frame registration
accuracy is comparable and long-horizon drift (ATE) is the discriminator.
On HeLiPR with C1 enabled, IV-GICP reduces RPE-t by
29\% (DCC05), 19\% (KAIST05), and 40\% (RIVER04), demonstrating that
C1 improves not only trajectory-level drift but also individual
frame alignment quality.

% ─────────────────────────────────────────────────────────────────────────────
\section{Discussion}\label{sec:discussion}

\textbf{Why two separate contributions rather than a monolithic system?}
C1 and C2 address \emph{different failure modes}: C1 re-distributes
optimization effort among existing geometric correspondences;
C2 injects an independent information channel.
They are complementary in principle but interact destructively on
uniform surfaces (Sec.~\ref{sec:c2p}); C2$'$ resolves this interaction
when both are needed.
In practice, the optimal configuration is determined by the surface
diversity: uniform tunnels use C1 alone ($\alpha{=}0$),
mixed-material environments use C2 alone ($\alpha{=}0.1$),
and C2$'$ bridges the two regimes.

\textbf{Zero tunable parameters.}
Via C1--$\alpha$ coupling, IV-GICP's two environment-dependent
choices (alpha and C1 on/off) reduce to a \emph{single binary decision}
determined by a measurable physical quantity:
the per-voxel intensity variance $\bar{V}_I$ of the first $N$ frames.
When $\bar{V}_I < \tau$ (uniform surfaces: concrete, Ouster
reflectivity), the system selects $\alpha{=}0$ with C1 on;
when $\bar{V}_I \geq \tau$ (diverse materials), $\alpha{=}0.1$ with
C1 off.
This is not tuning but \emph{data-driven measurement}: the same
physical quantity (intensity variance) that determines
$\omega_I^{(j)}$ in~\eqref{eq:4d_omega} also determines the
operating mode.
The remaining parameters (voxel size, max range, map radius) are
sensor- and range-determined, identical in nature to KISS-ICP's
per-environment defaults~\cite{vizzo2023kiss}.
IV-GICP thus has \textbf{zero free parameters} beyond sensor
specifications --- all choices derive from measurable quantities
or theoretical guarantees (Corollary~\ref{cor:alpha},
Theorem~\ref{thm:degeneracy}).

\textbf{Open-area limitation.}
Distribution-based methods assume sufficient per-voxel point density
for reliable covariance estimation.
With C1 enabled, most open-area gaps are resolved: RIVER04 improves
from $+15.9\%$ to $\mathbf{-33.2\%}$, DCC05 from $+21.6\%$ to
$-4.0\%$.  Only DCC04 (open convention center) retains a residual
$+15.5\%$ gap.

\begin{figure}[t]
  \centering
  \includegraphics[width=\columnwidth]{figures/c1_effect_bar.pdf}
  \caption{C1 (FIM weighting) ATE change by environment.
    Negative (blue) = improvement. C1 is essential for degenerate
    tunnels ($-23\%$) but harmful in well-conditioned outdoor ($+79\%$).}
  \label{fig:c1_effect}
\end{figure}

\textbf{C1--$\alpha$ coupling.}
C1 activation is principally coupled to the intensity weight $\alpha$
rather than requiring a separate environment detector.
When $\alpha{=}0$ (geometry-only), the total FIM $\mathcal{I}_{\mathrm{total}}
= \mathcal{I}_G$ may be rank-deficient along degenerate directions
(Theorem~\ref{thm:degeneracy}), and C1's per-correspondence FIM
reweighting is the sole defense.
When $\alpha{>}0$, Theorem~\ref{thm:degeneracy} guarantees
$\mathcal{I}_{\mathrm{total}} \succ 0$ regardless of geometric
conditioning; C1's aggressive reweighting then over-corrects and
disrupts the FIM balance ($+79\%$ on KITTI with $\alpha{=}0.1$).
This yields a two-mode selection rule:
\begin{itemize}
  \item $\alpha{=}0$, C1 on: uniform/degenerate surfaces (concrete tunnels,
    Ouster reflectivity noise) --- C1 provides the missing FIM regularization.
  \item $\alpha{>}0$, C1 off: diverse-material outdoor scenes (KITTI, SubT)
    --- Theorem~\ref{thm:degeneracy} already ensures well-posedness.
\end{itemize}
This mirrors the standard practice of per-environment parameter selection
(\eg KISS-ICP's voxel-size defaults~\cite{vizzo2023kiss}), but is
grounded in the FIM well-posedness condition rather than ad-hoc tuning.
We investigated automatic per-frame gating via the Hessian condition
number $\kappa_H$ and voxel-normal scatter $\kappa_N$, but both metrics
converge as the map grows ($\kappa_H{>}10^6$ for both tunnel and outdoor
at 500 frames), making threshold-based auto-gating infeasible at long
horizons.
A lighter spectral-gap gate on the 6-DOF FIM,
$g = (\lambda_2{-}\lambda_1)/\lambda_2 < \tau$, is nevertheless safe to
apply on outdoor data and does not regress KITTI at full horizon
(Appendix~\ref{app:fullseq}); inertial fusion remains the principled
long-horizon solution for kilometer-scale tunnels%
~\cite{pfreundschuh2024coinlio}.

\textbf{Comparison with IMU-based methods.}
COIN-LIO~\cite{pfreundschuh2024coinlio} achieves excellent tunnel
odometry (0.49--0.74\,m on ENWIDE tunnels) but requires tightly-coupled
IMU integration; without it, all tested methods --- including
FAST-LIO2~\cite{xu2022fastlio2} and LIO-SAM~\cite{shan2020liosam}
--- diverge in the same tunnels.
IV-GICP demonstrates that \emph{comparable degeneracy robustness is
achievable without inertial sensing}, via the FIM framework (C1+C2).
This is particularly relevant for survey platforms (push-carts,
handheld scanners) where IMU calibration is difficult or unavailable.

\textbf{Other limitations.}
(i)~KITTI processing speed ($4.7\times$ slower than KISS-ICP) is
dominated by multi-iteration GN on unbounded maps; with
$R_{\mathrm{map}}$ enabled, IV-GICP runs $1.4$--$1.9\times$ faster
on dense sensors.
(ii)~LiDAR-only: no IMU integration limits robustness to aggressive
motion (\eg rapid rotation with insufficient scan overlap).

% ─────────────────────────────────────────────────────────────────────────────
\section{Conclusion}\label{sec:conclusion}

IV-GICP derives two complementary contributions from a unified
Fisher Information maximization framework:
C1 (FIM-weighted correspondences) eliminates geometric degeneracy
in tunnels and corridors by concentrating optimization on informative
directions;
C2 (entropy-regularized 4D registration) provides principled intensity
integration with provable well-posedness (Theorem~\ref{thm:degeneracy}).
C1 activation is principally coupled to the intensity weight $\alpha$:
when $\alpha{=}0$ (geometry-only), C1 compensates for potential FIM
rank deficiency; when $\alpha{>}0$, Theorem~\ref{thm:degeneracy}
already ensures well-posedness, making C1 unnecessary.
On the challenging GEODE Urban Tunnel benchmark, IV-GICP reduces ATE
by 47--86\% vs.\ KISS-ICP and 12--83\% vs.\ GenZ-ICP while running
$1.8\times$ faster.
Standard VGICP catastrophically fails in the same conditions
(16--206\,m ATE), confirming the necessity of FIM-based weighting.
Across 7 datasets spanning 5 LiDAR sensors and 31 sequences, IV-GICP
achieves an 81\% win rate (25/31) vs.\ KISS-ICP, including $-24\%$
average improvement on HeLiPR Ouster data with C1.
On long-horizon evaluation, IV-GICP (i)~leads the KITTI-standard drift
metric on 10 of 11 full odometry sequences
(Table~\ref{tab:kitti-fullseq}, avg $0.867\%$ vs.\ $1.053\%$ for KISS
and $1.033\%$ for GenZ), and (ii)~beats KISS-ICP by $-42\%$ on the full
10{,}810-frame HeLiPR DCC05 sequence
(Table~\ref{tab:helipr-fullseq}), confirming that the C1/C2 framework
scales beyond the 500\,fr HD-map tile to multi-kilometer operation
in both diverse-material (KITTI) and dense-Ouster (HeLiPR) regimes.

Future work will extend the FIM framework to
factor-graph SLAM for large-scale HD map construction
in GPS-denied underground environments, and investigate
learned per-frame C1 gating to replace the current
dataset-level selection.

% ═════════════════════════════════════════════════════════════════════════════
% APPENDIX
% ═════════════════════════════════════════════════════════════════════════════
\appendix

\section{Mathematical Background}\label{app:math}

This appendix provides detailed derivations for the theoretical
results summarized in the main text.

\subsection{GICP Formulation}\label{app:gicp}

Generalized ICP~\cite{segal2009generalized} models both source and target
as local Gaussian distributions. Given source point $p_s$ with covariance
$C_s$ and target distribution $\mathcal{N}(\mu_t, C_t)$, the transformed
residual under pose $T = (R, t) \in \SE$ is:
\begin{equation}
  d(T) = Tp_s - \mu_t = Rp_s + t - \mu_t.
  \label{eq:app_residual}
\end{equation}
The combined distribution-to-distribution distance uses the Mahalanobis norm:
\begin{equation}
  D^2(T) = d(T)^T \bigl(C_t + R\,C_s\,R^T\bigr)^{-1} d(T).
  \label{eq:app_mahal}
\end{equation}
Summing over $M$ correspondences yields the negative log-likelihood
objective in~\eqref{eq:loglik}.

\textbf{Source covariance.}
We adopt the voxelized source covariance from VGICP~\cite{koide2021voxel}:
each source point inherits the sample covariance of its voxel,
regularized with $\sigma_{\min}^2 I_3$ to prevent singularity.

\textbf{Target covariance.}
Target voxels maintain online Welford statistics (Appendix~\ref{app:welford});
the resulting sample covariance is regularized as:
\begin{equation}
  \Sigma_j = \frac{1}{n_j - 1}\sum_{i=1}^{n_j}(x_i - \mu_j)(x_i - \mu_j)^T
  + \sigma_{\min}^2 I_3.
  \label{eq:app_target_cov}
\end{equation}

\subsection{Fisher Information Matrix Derivation}\label{app:fim}

The pose is parameterized as $T(\xi) = \exp(\xi^\wedge)\,T_0$
with $\xi = [\omega^T, \rho^T]^T \in \mathfrak{se}(3)$ (rotation then
translation).  The Jacobian of the residual with respect to $\xi$ is:
\begin{equation}
  J_m = \frac{\partial d_m}{\partial \xi}\bigg|_{\xi=0}
  = \begin{bmatrix} -(Rp_m)^\times & I_3 \end{bmatrix} \in \mathbb{R}^{3\times 6},
  \label{eq:app_jacobian}
\end{equation}
where $(\cdot)^\times$ denotes the skew-symmetric matrix.

The Fisher Information Matrix (FIM) for the full pose is:
\begin{equation}
  \mathcal{I}(\xi) = \sum_{m=1}^{M} J_m^T \,\Omega_m\, J_m
  \in \mathbb{R}^{6\times 6},
  \label{eq:app_fim}
\end{equation}
with precision $\Omega_m = (C_t^{(j)} + R\,C_s^{(m)}\,R^T)^{-1}$.

The Gauss-Newton update solves:
\begin{equation}
  \mathcal{I}(\xi)\,\delta\xi = -\sum_{m} J_m^T \Omega_m\, d_m,
  \label{eq:app_gn}
\end{equation}
and applies $T \leftarrow \exp(\delta\xi^\wedge)\,T$.

\textbf{Eigendecomposition.}
Let $\lambda_1 \geq \cdots \geq \lambda_6$ be eigenvalues of $\mathcal{I}$.
The condition number $\kappa = \lambda_1/\lambda_6$ quantifies degeneracy.
When $\lambda_6 \approx 0$, the GN step along $v_6$ is numerically unstable.

\subsection{Theorem 1: Degeneracy Recovery (Full Proof)}\label{app:thm1}

\begin{theorem*}[Restated]
If $\alpha > 0$ and at least one voxel has $\|\nabla\mu_j^I\|_2 > 0$,
then $\mathcal{I}_{\mathrm{total}} \succ 0$.
\end{theorem*}

\begin{proof}
Decompose the total FIM:
\begin{equation}
  \mathcal{I}_{\mathrm{total}} = \underbrace{\sum_m J_m^{3\times6,T}
  \Omega_m^{\mathrm{geo}} J_m^{3\times6}}_{\mathcal{I}_G}
  + \underbrace{\sum_m (j_m^I)^T \omega_I^{(m)} j_m^I}_{\mathcal{I}_I},
  \label{eq:app_fim_decomp}
\end{equation}
where $j_m^I = \alpha\,(\nabla\mu_m^I)^T J_m^{xyz} \in \mathbb{R}^{1\times 6}$
is the intensity Jacobian row for correspondence~$m$.

For any unit vector $v \in S^5$:
\begin{align}
  v^T \mathcal{I}_{\mathrm{total}}\, v
  &= v^T \mathcal{I}_G\, v + v^T \mathcal{I}_I\, v \nonumber\\
  &\geq v^T \mathcal{I}_I\, v
  = \sum_m \omega_I^{(m)} (j_m^I\, v)^2 \geq 0.
  \label{eq:app_bound}
\end{align}
It remains to show the intensity term is strictly positive for all $v$.

Form the $M \times 6$ matrix $\mathbf{J}_I = [j_1^I;\, \ldots;\, j_M^I]$.
Then $\mathcal{I}_I = \mathbf{J}_I^T\, \mathrm{diag}(\omega_I)\,\mathbf{J}_I$.
The null space of $\mathcal{I}_I$ equals that of
$\mathrm{diag}(\omega_I^{1/2})\,\mathbf{J}_I$.

Each row $j_m^I = \alpha\,(\nabla\mu_m^I)^T\, [-(Rp_m)^\times\;\; I_3]$.
If at least 6 correspondences have non-collinear intensity gradients
(which is always the case with $M > 1000$ real-world correspondences),
then $\mathrm{rank}(\mathbf{J}_I) = 6$, and
$v^T\mathcal{I}_I\,v > 0$ for all $v \neq 0$.

Hence $\mathcal{I}_{\mathrm{total}} = \mathcal{I}_G + \mathcal{I}_I \succ 0$.
\end{proof}

\textbf{Practical implication.}
When $\alpha > 0$, the intensity channel provides a ``safety net'' for
any degenerate direction: even if $\mathcal{I}_G$ has a zero eigenvalue
(\eg tunnel axis), $\mathcal{I}_I$ fills the gap.
This is why C1 (FIM reweighting) is unnecessary when $\alpha > 0$ ---
the well-posedness is already guaranteed.

\subsection{Theorem 2: Propagation Error Bound (Full Proof)}\label{app:thm2}

\begin{theorem*}[Restated]
With odometry uncertainty $\Sigma_{\Delta T}$:
\begin{equation}
  \|\Sigma_j^{\mathrm{true}} - R_\Delta\Sigma_j^{(k)}R_\Delta^T\|_F
  \leq \|\mu_j^{(k)}\|^2 \|\Sigma_{\Delta T}\|_F + O(\|\Sigma_{\Delta T}\|_F^2).
\end{equation}
\end{theorem*}

\begin{proof}
Let $T_{\mathrm{true}} = \hat{T}\,\exp(\delta\xi^\wedge)$ where
$\delta\xi \sim \mathcal{N}(0, \Sigma_{\Delta T})$.
A point $x$ in the voxel frame transforms as:
\begin{align}
  T_{\mathrm{true}}\,x &= \hat{T}\,\exp(\delta\xi^\wedge)\,x \nonumber\\
  &= \hat{T}\,(I + \delta\xi^\wedge + O(\|\delta\xi\|^2))\,x \nonumber\\
  &= \hat{T}\,x + \hat{T}\,\delta\xi^\wedge\,x + O(\|\delta\xi\|^2).
\end{align}
The transformed covariance is:
\begin{equation}
  \Sigma^{\mathrm{true}} = R_\Delta\,\Sigma\,R_\Delta^T
  + R_\Delta\,\mathbb{E}[\delta\xi^\wedge x x^T
  (\delta\xi^\wedge)^T]\,R_\Delta^T.
\end{equation}
The perturbation term involves:
\begin{align}
  \|\delta\xi^\wedge x\|^2 &= \|\delta\omega \times x + \delta\rho\|^2 \nonumber\\
  &\leq (\|\delta\omega\|\|x\| + \|\delta\rho\|)^2 \nonumber\\
  &\leq 2(\|x\|^2\|\delta\omega\|^2 + \|\delta\rho\|^2).
\end{align}
Taking expectation and using $\mathbb{E}[\|\delta\xi\|^2] = \mathrm{tr}(\Sigma_{\Delta T})
\leq 6\|\Sigma_{\Delta T}\|_F$:
\begin{equation}
  \|\Sigma^{\mathrm{true}} - R_\Delta\Sigma R_\Delta^T\|_F
  \leq C\,\|x\|^2\,\|\Sigma_{\Delta T}\|_F,
\end{equation}
where $\|x\| \leq \|\mu_j^{(k)}\| + O(\sqrt{\mathrm{tr}(\Sigma_j)})$.
For well-localized voxels ($\mathrm{tr}(\Sigma_j)$ small),
$\|x\| \approx \|\mu_j^{(k)}\|$.
\end{proof}

\textbf{Consequence.}
As odometry improves ($\Sigma_{\Delta T} \to 0$), map distribution error
vanishes, creating a virtuous cycle: better odometry $\to$ better map $\to$
better correspondences $\to$ better odometry.

\subsection{Welford Online Update}\label{app:welford}

The VoxelMap maintains per-voxel statistics using Welford's online
algorithm~\cite{welford1962note}, which computes running mean and
covariance in $O(1)$ per insertion without storing individual points.

For voxel $j$ with $n$ points inserted:
\begin{align}
  \mu_j^{(n)} &= \mu_j^{(n-1)} + \frac{x_n - \mu_j^{(n-1)}}{n}, \\
  M_2^{(n)} &= M_2^{(n-1)} + (x_n - \mu_j^{(n-1)})(x_n - \mu_j^{(n)})^T, \\
  \Sigma_j^{(n)} &= \frac{M_2^{(n)}}{n-1} + \sigma_{\min}^2 I_3.
\end{align}
The intensity mean $\mu_j^I$ and variance $\mathrm{Var}_j(I)$ are
similarly maintained via scalar Welford updates.

\textbf{Numerical stability.}
Welford's algorithm is numerically stable for large $n$ because it
avoids the catastrophic cancellation of the naive two-pass formula
$\Sigma = \frac{1}{n-1}(\sum x_i x_i^T - n\bar{x}\bar{x}^T)$.
This is critical for long-running odometry where voxels accumulate
thousands of points.

\subsection{C1 Weight Derivation}\label{app:c1}

\textbf{Goal}: allocate per-correspondence weights $\{w_m\}$ to
maximize $\lambda_{\min}(\mathcal{I})$, \ie strengthen the weakest
direction of the FIM.

The weighted FIM is:
\begin{equation}
  \mathcal{I}(w) = \sum_m w_m\, J_m^T \Omega_m\, J_m.
\end{equation}
The gradient of $\lambda_{\min}$ with respect to $w_m$ is:
\begin{equation}
  \frac{\partial \lambda_{\min}}{\partial w_m}
  = v_{\min}^T (J_m^T \Omega_m J_m)\, v_{\min}
  = w_m^{\mathrm{FIM}},
\end{equation}
where $v_{\min}$ is the eigenvector corresponding to $\lambda_{\min}$.

The greedy ascent direction allocates weight proportional to
$w_m^{\mathrm{FIM}}$: correspondences with larger FIM contribution
along $v_{\min}$ receive higher weight.

Combined with Huber robustness ($w_m^H$, which down-weights outliers):
\begin{equation}
  w_m = w_m^H \cdot \frac{w_m^{\mathrm{FIM}}}{\bar{w}^{\mathrm{FIM}} + \varepsilon_0},
\end{equation}
where $\bar{w}^{\mathrm{FIM}} = M^{-1}\sum_m w_m^{\mathrm{FIM}}$ normalizes
the scale and $\varepsilon_0$ prevents division by zero when
$\mathcal{I}$ is exactly rank-deficient.

\subsection{C2$'$: Per-Voxel Adaptive Gate}\label{app:c2p}

When C1 and C2 are simultaneously active, a destructive interaction occurs:
C1 suppresses degenerate-direction correspondences while C2 assigns high
intensity precision $\omega_I$ to the same correspondences (uniform
concrete $\to$ low $\mathrm{Var}(I)$ $\to$ high $\omega_I$).

The gate uses the geometric condition number of the \emph{unregularized}
voxel covariance $\hat{\Sigma}_j = M_2^{(n)}/n + \varepsilon I_3$
(with $\varepsilon = 10^{-9}$, not the sigma-floor $\sigma_{\min}^2$):
\begin{equation}
  \kappa_j = \frac{\lambda_{\max}(\hat{\Sigma}_j)}{\lambda_{\min}(\hat{\Sigma}_j)}.
\end{equation}

The sigmoid gate:
\begin{equation}
  g_j = \sigma\!\bigl(\log(\kappa_j / \kappa_0)\bigr)
  = \frac{1}{1 + (\kappa_0/\kappa_j)},
\end{equation}
where $\kappa_0 = 50$ is the transition threshold.

The gated intensity precision becomes:
\begin{equation}
  \omega_I^{(j)} \leftarrow g_j \cdot \omega_I^{(j)}.
\end{equation}

\textbf{Behavior:}
\begin{itemize}
  \item $\kappa_j \gg \kappa_0$ (degenerate, \eg tunnel wall):
    $g_j \to 1$, full intensity precision.
  \item $\kappa_j \ll \kappa_0$ (well-conditioned):
    $g_j \to 0$, geometry-only.
\end{itemize}

This resolves the C1+C2 destructive interference:
on GEODE Urban Tunnel01, ungated C1+C2 gives ATE = 1.50\,m ($+52\%$
vs base), while C2$'$ reduces this to 0.85\,m ($-14\%$ vs base).

\subsection{4D Precision Matrix Structure}\label{app:4d}

The full 4D precision for correspondence $m$ with target voxel $j$ is:
\begin{equation}
  \Omega_m^{4\times4} =
  \begin{bmatrix}
    (\Sigma_j^{3\times3} + R\Sigma_s R^T)^{-1} & \mathbf{0} \\
    \mathbf{0}^T & \omega_I^{(j)}
  \end{bmatrix},
\end{equation}
where the block-diagonal structure reflects conditional independence
of position and intensity given voxel assignment.

The intensity precision is:
\begin{equation}
  \omega_I^{(j)} = \frac{\alpha^2}{\mathrm{Var}_j(I)/\ell_v^2 + \varepsilon_I},
\end{equation}
where $\ell_v$ is the voxel size, $\mathrm{Var}_j(I)$ is the sample
intensity variance from Welford updates, and $\varepsilon_I$ prevents
infinite precision for zero-variance voxels.

The 4D Jacobian is:
\begin{equation}
  J_m^{4\times6} =
  \begin{bmatrix}
    -(Rp_m)^\times & I_3 \\
    -\alpha\,(\nabla\mu_j^I)^T (Rp_m)^\times & \alpha\,(\nabla\mu_j^I)^T
  \end{bmatrix}.
\end{equation}

\subsection{Convergence and Initialization}\label{app:convergence}

The Gauss-Newton iteration converges when $T_0$ is within the basin
of attraction. Constant-velocity initialization:
\begin{equation}
  T_0 = T_{k-1} \cdot T_{k-2}^{-1} \cdot T_{k-1}.
\end{equation}

The adaptive correspondence distance:
\begin{equation}
  \sigma_k = \max\!\bigl(\sigma_{\min},\;
  \mathrm{median}\bigl\{\|T_k p_m - \mu_{j(m)}\|\bigr\}\bigr),
\end{equation}
where $\sigma_{\min}$ (0.1--0.5\,m) prevents cascade collapse in
degenerate geometry.

\subsection{Full-Sequence Verification}\label{app:fullseq}

To verify that the 500\,fr evaluation window of Section~\ref{sec:experiments}
does not hide long-horizon regressions, we ran full-sequence benchmarks on
(a)~KITTI odometry sequences 00--10, using the KITTI-standard relative
translational drift [\%] over $\{100,\ldots,800\}$\,m
segments~\cite{geiger2013kitti} as reported by
KISS-ICP~\cite{vizzo2023kiss} and GenZ-ICP~\cite{zhang2024genz}, and
(b)~GEODE Urban Tunnels 01--03 (2857--3425\,fr, $\sim$3\,km of
GNSS-denied driving), reporting ATE as in GenZ-ICP's degenerate-scene
evaluation.

\textbf{KITTI full-sequence drift (\%).}
Table~\ref{tab:kitti-fullseq} reports the KITTI-standard drift metric
on all 11 sequences. IV-GICP achieves the lowest average drift
($0.867\%$), beating KISS-ICP ($1.053\%$, $-17.7\%$) and
GenZ-ICP ($1.033\%$, $-16.1\%$), and wins the drift metric on
$10/11$ sequences. The per-sequence average ATE is
$4.09$\,m (IV), $3.91$\,m (KISS), and $3.45$\,m (GenZ), within normal
long-horizon variance. The absolute drift numbers differ from the
authors' originally reported averages
(KISS $0.50\%$~\cite{vizzo2023kiss}, GenZ $0.51\%$~\cite{zhang2024genz});
we attribute this to our single-machine re-runs without their full
deskewing/IMU-free re-tuning, and note that the \emph{relative}
ordering and separations are the comparable quantity.

\begin{table}[h]
\centering
\scriptsize
\setlength{\tabcolsep}{3pt}
\caption{KITTI full-sequence 3-way comparison. Drift\,[\%] = mean
relative translational error on 100--800\,m segments (KITTI benchmark
standard). ATE = Umeyama-aligned RMSE [m]. Bold = best per row and
metric.}
\label{tab:kitti-fullseq}
\begin{tabular}{l r | r r r | r r r}
\toprule
 & & \multicolumn{3}{c|}{Drift\,[\%]} & \multicolumn{3}{c}{ATE\,[m]} \\
Seq & n & KISS & IV & GenZ & KISS & IV & GenZ \\
\midrule
00 & 4541 & 0.906 & \textbf{0.694} & 0.924 & 3.995 & 5.563 & \textbf{2.519} \\
01 & 1101 & 2.120 & \textbf{1.829} & 2.238 & 19.822 & \textbf{17.205} & 20.282 \\
02 & 4661 & 1.227 & \textbf{1.054} & 1.133 & 7.404 & 10.625 & \textbf{4.771} \\
03 &  801 & 1.071 & 1.207 & \textbf{1.050} & 0.852 & \textbf{0.847} & 0.867 \\
04 &  271 & 0.928 & \textbf{0.825} & 0.935 & 0.420 & \textbf{0.286} & 0.435 \\
05 & 2761 & 0.778 & \textbf{0.617} & 0.666 & \textbf{1.738} & 3.037 & 1.137 \\
06 & 1101 & 0.642 & \textbf{0.541} & 0.666 & 0.896 & \textbf{0.658} & 0.842 \\
07 & 1101 & 0.495 & \textbf{0.486} & 0.522 & 0.355 & 0.624 & \textbf{0.348} \\
08 & 4071 & 1.112 & \textbf{0.921} & 1.035 & 4.079 & 3.539 & \textbf{3.492} \\
09 & 1591 & 0.899 & \textbf{0.584} & 0.984 & 1.585 & 1.668 & \textbf{1.523} \\
10 & 1201 & 1.404 & \textbf{0.782} & 1.214 & 1.880 & \textbf{0.937} & 1.711 \\
\midrule
avg & & 1.053 & \textbf{0.867} & 1.033 & 3.911 & 4.090 & \textbf{3.448} \\
\bottomrule
\end{tabular}
\end{table}

\textbf{GEODE Urban Tunnel full-sequence ATE (m).}
On the complete 3\,km tunnels, every tested LiDAR-only method diverges
catastrophically: KISS-ICP 716--1092\,m, IV-GICP 691--1524\,m, and
GenZ-ICP 479--713\,m ATE (full table in the supplementary material).
This reproduces COIN-LIO's finding that LiDAR-only odometry (including
KISS-ICP, FAST-LIO2, and LIO-SAM without inertial aid) fails on
multi-kilometer tunnels~\cite{pfreundschuh2024coinlio}, and confirms
that meaningful comparison in degenerate environments requires an
evaluation window shorter than the drift-explosion horizon. Our 500\,fr
window ($\approx$50\,s at 10\,Hz) sits safely inside this operational
regime for every dataset considered.

\textbf{HeLiPR DCC05 long-horizon ATE (10{,}810\,fr).}
To verify long-horizon robustness outside a diverging-tunnel regime, we
ran the full HeLiPR DCC05 sequence (Ouster OS1-64, 10{,}810 frames,
campus roads with sharp turns).  Enforcing the sigma floor
$\sigma_{\min}{=}0.5$\,m (Sec.~\ref{sec:subt}) --- which prevents the
adaptive correspondence radius from collapsing during aggressive
rotations --- yields the results in
Table~\ref{tab:helipr-fullseq}.
IV-GICP achieves $17.21$\,m ATE against KISS-ICP's $29.67$\,m
($\mathbf{-42.0\%}$), a $22{\times}$ improvement over the same pipeline
with $\sigma_{\min}{=}0.1$\,m.  Together with the KITTI drift-\%
result (Table~\ref{tab:kitti-fullseq}), this provides \emph{two
independent long-horizon datasets} --- 4.1\,km of KITTI driving and
10.8\,k frames of HeLiPR campus --- on which the C1/C2 framework
maintains (KITTI) or \emph{improves upon} (HeLiPR) the KISS-ICP/GenZ-ICP
baselines at full sequence length.

\begin{table}[h]
\centering
\small
\caption{HeLiPR DCC05 full-sequence ATE [m] (10{,}810\,fr, Ouster OS1-64).
The sigma floor $\sigma_{\min}{=}0.5$\,m is the
same mechanism used for the SubT-MRS results of Table~\ref{tab:subt}.}
\label{tab:helipr-fullseq}
\begin{tabular}{lcc}
\toprule
Method & ATE [m] & vs.\ KISS \\
\midrule
IV-GICP, $\sigma_{\min}{=}0.1$\,m (no floor) & 378.44 & $+1175\%$ \\
KISS-ICP                                      & 29.67  & baseline \\
\textbf{IV-GICP, $\sigma_{\min}{=}0.5$\,m (floor)} & \textbf{17.21} & $\mathbf{-42.0\%}$ \\
\bottomrule
\end{tabular}
\end{table}

\textbf{Auto-gate non-regression.}
A gap-ratio gate on the 6-DOF geometric FIM eigenvalues,
$g = (\lambda_2{-}\lambda_1)/\lambda_2 < \tau$, with
$\tau\in\{0.005,0.02,0.05\}$, was evaluated on KITTI seq\,00
full-length (4541\,fr). All three thresholds yield ATE
$4.73$--$5.09$\,m, within $\pm7\%$ of C1-off (5.09\,m), whereas
manual C1-on regresses to $7.45$\,m ($+46\%$). The gate is therefore
safe on outdoor data while recovering the tunnel benefits of manual C1
at short horizons. It does not help at full tunnel horizons because
the eigenvalue gap stabilizes as the voxel map accumulates wall
observations, the same saturation noted by prior degeneracy-analysis
work~\cite{zhang2016degeneracy}.

% ─────────────────────────────────────────────────────────────────────────────
\bibliographystyle{IEEEtran}
\begin{thebibliography}{99}

\bibitem{vizzo2023kiss}
I.~Vizzo, T.~Guadagnino, B.~Mersch, L.~Nardi, J.~Behley, and C.~Stachniss,
``KISS-ICP: In Defense of Point-to-Point ICP -- Simple, Accurate, and Robust
Registration If Done the Right Way,''
\textit{IEEE RA-L}, vol.~8, no.~2, pp.~1029--1036, 2023.

\bibitem{koide2021voxel}
K.~Koide, M.~Yokozuka, S.~Oishi, and A.~Banno,
``Voxelized GICP for Fast and Accurate 3D Point Cloud Registration,''
in \textit{Proc. IEEE ICRA}, 2021, pp.~11054--11059.

\bibitem{zhang2024genz}
J.~Zhang, P.~Pfreundschuh, and R.~Siegwart,
``GenZ-ICP: Generalizable and Degeneracy-Robust LiDAR Odometry Using an
Adaptive Weighting,''
\textit{IEEE RA-L}, 2025. [arXiv:2411.06766]

\bibitem{pfreundschuh2024coinlio}
P.~Pfreundschuh, H.~Oleynikova, C.~Cadena, R.~Siegwart, and O.~Andersson,
``COIN-LIO: Complementary Intensity-Augmented LiDAR Inertial Odometry,''
in \textit{Proc. IEEE ICRA}, 2024. [arXiv:2310.01235]

\bibitem{geiger2013kitti}
A.~Geiger, P.~Lenz, C.~Stiller, and R.~Urtasun,
``Vision Meets Robotics: The KITTI Dataset,''
\textit{Int. J. Rob. Res.}, vol.~32, no.~11, pp.~1231--1237, 2013.

\bibitem{zhang2014loam}
J.~Zhang and S.~Singh,
``LOAM: Lidar Odometry and Mapping in Real-time,''
in \textit{Proc. RSS}, 2014.

\bibitem{shan2018lego}
T.~Shan and B.~Englot,
``LeGO-LOAM: Lightweight and Ground-Optimized Lidar Odometry and Mapping
on Variable Terrain,''
in \textit{Proc. IEEE/RSJ IROS}, 2018, pp.~4758--4765.

\bibitem{vizzo2024degeneracy}
I.~Vizzo \etal,
``Degeneracy-Aware Point Cloud Registration,''
arXiv:2408.11809, 2024.

\bibitem{shan2021geode}
T.~Shan, B.~Englot, D.~Meyers, W.~Wang, C.~Ratti, and D.~Rus,
``GEODE: Generalized Environmental Object Detection via 3D LiDAR Mapping,''
arXiv:2104.04950, 2021.

\bibitem{umeyama1991least}
S.~Umeyama,
``Least-Squares Estimation of Transformation Parameters Between Two Point
Patterns,''
\textit{IEEE Trans.\ PAMI}, vol.~13, no.~4, pp.~376--380, 1991.

\bibitem{segal2009generalized}
A.~Segal, D.~H\"{a}hnel, and S.~Thrun,
``Generalized-ICP,''
in \textit{Proc.\ RSS}, 2009.

\bibitem{helmberger2022hilti}
D.~Helmberger, K.~Morin, B.~Berner, N.~Kumar, G.~Cioffi, and D.~Scaramuzza,
``The Hilti SLAM Challenge Dataset,''
\textit{IEEE RA-L}, vol.~7, no.~3, pp.~7518--7525, 2022.

\bibitem{subt2021}
T.~Rogers \etal,
``SubT-MRS: Pushing SLAM towards All-weather Environments,''
arXiv:2307.07607, 2023.

\bibitem{jung2023helipr}
M.~Jung \etal,
``HeLiPR: Heterogeneous LiDAR Dataset for inter-LiDAR Place Recognition
under Spatial and Temporal Variations,''
\textit{Int. J. Rob. Res.}, 2024.

\bibitem{kim2020mulran}
G.~Kim, Y.~S.~Park, Y.~Cho, J.~Jeong, and A.~Kim,
``MulRan: Multimodal Range Dataset for Urban Place Recognition,''
in \textit{Proc.\ IEEE ICRA}, 2020, pp.~6246--6253.

\bibitem{dellenbach2022ct}
J.~Dellenbach, J.-E.~Deschaud, B.~Jacquet, and F.~Goulette,
``CT-ICP: Real-time Elastic LiDAR Odometry with Loop Closure,''
in \textit{Proc.\ IEEE ICRA}, 2022, pp.~5580--5586.

\bibitem{chen2022dlo}
K.~Chen, B.~T.~Lopez, A.-A.~Agha-mohammadi, and A.~Mehta,
``Direct LiDAR Odometry: Fast Localization with Dense Point Clouds,''
\textit{IEEE RA-L}, vol.~7, no.~2, pp.~2000--2007, 2022.

\bibitem{xu2022fastlio2}
W.~Xu, Y.~Cai, D.~He, J.~Lin, and F.~Zhang,
``FAST-LIO2: Fast Direct LiDAR-Inertial Odometry,''
\textit{IEEE Trans.\ Rob.}, vol.~38, no.~4, pp.~2053--2073, 2022.

\bibitem{zhang2016degeneracy}
J.~Zhang, M.~Kaess, and S.~Singh,
``On Degeneracy of Optimization-based State Estimation Problems,''
in \textit{Proc.\ IEEE ICRA}, 2016, pp.~809--816.

\bibitem{shan2020liosam}
T.~Shan, B.~Englot, D.~Meyers, W.~Wang, C.~Ratti, and D.~Rus,
``LIO-SAM: Tightly-Coupled Lidar Inertial Odometry via Smoothing and
Mapping,''
in \textit{Proc.\ IEEE/RSJ IROS}, 2020, pp.~5135--5142.

\bibitem{welford1962note}
B.~P.~Welford,
``Note on a Method for Calculating Corrected Sums of Squares and Products,''
\textit{Technometrics}, vol.~4, no.~3, pp.~419--420, 1962.

\bibitem{potokar2025form}
E.~Potokar, S.~Olsen, J.~Bown, and J.~Mangelson,
``FORM: A Framework for Online Registration-based Mapping,''
in \textit{Proc.\ IEEE ICRA}, 2025. [arXiv:2510.09966]

\end{thebibliography}

\end{document}
